\subsection{问题一：并行粒度与Task边界}
问题一须在并行粒度与搬运代价之间取舍。切出更多Task可分散计算，却增加边界COPY、激活等待，以及M/V流水线和60~B/cycle DDR带宽上的排队。相同张量在同一接收Task内只读入一次，分散到多个Task才重复搬运；切分数与Makespan收益并非一一对应。

弱连通分量提供不切断依赖的划分起点；单一大分量则需沿拓扑序切分。切分依据计算周期与边界张量字节数，分核时还需考虑两条计算流水线的负载。周期、搬运和等待代价用于筛选候选顺序，Makespan由完整事件模拟给出。

\subsection{问题二：同核复用与容量重叠}
问题二将切分搬运的接收域由Task改为核心。同核合并可减少重复读入，却使更多物理张量版本在L1、UB中重叠存活，可能触发缓存换出换入（Spill）。因此，方案评价须同时核算切分搬运与Spill。

对于固定方案$X$，A、B的接收域与跨核写回规则不同。问题二按物理副本从申请到最后消费的区间检查L1、UB占用，迁核会同时改变接收域、核内顺序及存活区间的重叠。两问分别构造方案，以各自评估程序给出的Makespan和搬运量选优；容量压力估计仅用于筛选。

\subsection{问题三：查询时序与服务池}
问题三的Cache在读发射时查询、读完成时插入，并按FIFO淘汰。多个消费者的命中取决于查询与首次读完成的时序：后一次查询早于首次读完成时仍访问DDR；命中读还要竞争250~B/cycle的Cache服务池。命中率上升时，节省的读时延只有落在工期关键链上才会改变完成时刻。

同一逻辑张量的消费者聚到少数接收核心可少发COPY，却压缩并行；分到更多接收核心则可能在首读完成前增加DDR未命中。同图、同核数下固定方案在B、C两种配置下分别评价，可以把Cache硬件的作用与调度策略的调整分开；进一步在C配置下调整分核和次序，衡量方案适配的独立贡献。
